% Sec. I Introduction (about 1 page).
% Paragraphs: (1) the real-world gap; (2) IoRT and the surveys: why now;
% (3) why indoors is hard on every layer; (4) no existing platform;
% (5) SCooP-iT and its three uses; (6) contributions; (7) seven questions
% and the roadmap.
\section{Introduction}
\label{sec:introduction}
\IEEEPARstart{T}{eams} of robots and fixed sensors that share what they
perceive can see beyond the occlusions of any one of them. They
cooperate through collaborative SLAM, which places every agent in one
shared map \cite{swarmslam,doorslam,kimeramulti}, and cooperative
perception, which fuses what the agents observe within it
\cite{where2comm,zhou2022gnnperception}. The methods for both are
developed on vehicle datasets and in simulation
\cite{opv2v,v2xsim,dairv2x,v2v4real}, where agents carry the same
sensors, ground truth comes from GNSS, and messages arrive on time. A
real indoor team differs on every layer at once. Its sensors fall
short of their nominal rates and differ from agent to agent; its
clocks drift; there is no GNSS to tell where an agent was; the
building occludes cameras and radio alike; and its agents share one
contended wireless network, so a neighbor's message may arrive after
the scene has changed. Data freshness, studied as the age of
information \cite{kaul2012aoi,yates2021aoi}, here has a concrete
deadline of one sensor frame.

The Internet of Robotic Things treats such robots as mobile nodes that
sense and share data over the same network as fixed devices
\cite{simoens2018iort}. Recent surveys of networked embodied
intelligence draw the consequence for evaluation: delivering a message
is only an intermediate step whose value depends on whether it arrives
in time for the task, so sensing, communication, and the task outcome
have to be measured together \cite{zhang2026mean}. They and surveys of
embodied multi-robot systems name the missing instrument: benchmarks
evaluate communication, robotics, or perception separately, most in
simulations that overlook real sensing and communication, and
integrated testbeds with real wireless dynamics, heterogeneous agents,
and multi-agent tasks, bridged to simulation, are an open problem
\cite{zhang2026mean,wen2026embodiedmrs}.

Building such an instrument indoors is hard on every layer. Each
agent's sensing must be on one timeline with every other's, with a
bound on the error; each agent must have a pose in one metric frame
at building scale, without motion capture or GNSS; each link must be
recorded per agent pair, with the channel beneath it; and, for the
instrument to be used by others, the capture must be replayable into a
new method and extensible into situations it did not record.

No existing platform provides this. IoT and wireless testbeds measure
the radio but run no scored task
\cite{adjih2015iotlab,breen2021powder,raychaudhuri2020cosmos,bonati2021colosseum}.
Wi-Fi channel datasets measure CSI on dedicated sounding links, with
positions from a survey instrument \cite{dichasus,espargos}.
Communication-aware robot systems run over real links that they
monitor but do not release, and score nothing against a reference
\cite{ginting2021chord,saboia2022achord,cladera2024mocha}.
Cooperative-perception platforms are outdoor and GNSS-referenced, and
reduce the measured link to a scalar delay or a live run
\cite{v2v4real,zimmer2024tumtrafv2x,v2xrealo,cooperscene}. Embodied
multi-agent environments are simulated, with generated worlds and an
ideal network \cite{lingjing}.

This paper presents SCooP-iT, the indoor testbed of SCooP (Sensing and
Communication for Cooperative Perception): a real-world testbed in
which a heterogeneous team is captured on one clock and in one frame,
sensing and communication alike, and an environment built from that
capture. It is meant to be used in three ways. The \emph{recordings}
are a test bench: any collaborative method is replayed on what the
agents sensed and against the link they had. \emph{SCooP-iT Gym} is a
closed-loop environment whose world, sensors, link, and failure
conditions are built from the recordings, in which agents are trained
where the recordings cannot go. And the \emph{testbed} itself closes
the loop: a policy trained in the Gym is deployed on the real agents,
and what happened is compared with what the Gym predicted. Each use is
validated against reality. The contributions are:
\begin{enumerate}
    \item \textbf{A heterogeneous indoor testbed with sensing and
    communication on one clock}: two mobile and three infrastructure
    agents on one access point, every pair's link recorded (passive
    state at \SI{5}{\hertz}, RTT at \SI{100}{\hertz}, goodput in
    \SI{4}{\second} windows, CSI at \SIrange{104}{141}{\hertz}),
    contention as a controlled knob, and measured clock offsets below
    \SI{1}{\milli\second} (Secs.~\ref{sec:design}, \ref{sec:alignment},
    \ref{sec:link}).
    \item \textbf{Building-scale ground truth without motion capture,
    GNSS, or surveyed anchors}: a self-built reference map and a
    fiducial chain for agents without LiDAR, validated by measurements
    that never enter it, \todo{cross-pass figure} and
    \SI{10}{\milli\metre} against held-out fiducials
    (Sec.~\ref{sec:gt}).
    \item \textbf{A replay validated against live runs}: the recorded
    per-pair, per-message link, and the other recorded conditions, fed
    into any collaborative method keyed to sensor time,
    \todo{validated against live execution to within one receiver
    frame} (Sec.~\ref{sec:replay}, E3).
    \item \textbf{SCooP-iT Gym}: a closed-loop multi-agent environment
    with a Gymnasium interface, whose building, sensors, motion, link,
    and failure conditions are built from the recordings and each
    checked against held-out real runs, with every failure traced to
    its cause, and a train--test--deploy protocol that ends on the real
    testbed (Sec.~\ref{sec:gym}).
    \item \textbf{What the real world shows that simulation hides}: with
    a second robot on the access point, the loaded round trip of
    \SIrange{25}{33}{\milli\second} stays well inside one
    \SI{100}{\milli\second} LiDAR frame, yet only
    \SIrange{59}{68}{\percent} of \SI{1.5}{\mega\byte} messages arrive
    within it (86\,\% with one robot), while the synthetic channels of
    prior benchmarks deliver 100\,\% (Sec.~\ref{sec:experiments}).
\end{enumerate}

A testbed paper is judged on seven questions, five about whether the
instrument can be trusted and two about whether it is worth having
(Table~\ref{tab:spec}): \emph{Is it precise?} Are its own errors small
next to the effects it measures (E1)? \emph{Is it consistent?} Do the
poses and the link logs agree when neither was checked by hand (E2)?
\emph{Is it faithful?} Do the replay and the Gym reproduce what the
real agents sensed, sent, and achieved (E3)? \emph{Is it repeatable?}
Does the same setting give the same link on another day (E4)? \emph{Is
it necessary?} What does the real world show that a synthetic one
hides (E5)? \emph{Is it unique?} What can it do that no existing
platform can (Sec.~\ref{sec:related})? \emph{Is it reproducible?} Can
someone else build and use it (Secs.~\ref{sec:design},
\ref{sec:lessons})? Each trust question has a pass criterion fixed
before the data were read. Sec.~\ref{sec:related} reviews related work;
Secs.~\ref{sec:design}--\ref{sec:replay} describe the testbed and its
replay; Sec.~\ref{sec:experiments} reports the experiments;
Sec.~\ref{sec:gym} presents the Gym and its train--test--deploy
protocol; and Sec.~\ref{sec:lessons} gives the lessons the real world
taught and the limitations. The sequences and labels recorded on the
testbed are released as the SCooP dataset \cite{scoop_dataset}.
